\section{Población y muestra}
\label{sec:poblacion_muestra}

\subsection{Población y muestra humana}

La población humana del estudio comprende a la totalidad de estudiantes del nivel secundario procedentes de las diversas unidades educativas del distrito municipal de San Julián que aspiran a cursar el Bachillerato Técnico Humanístico en el MÓDULO TECNOLÓGICO PRODUCTIVO SAN JULIÁN, estimada en aproximadamente 400 postulantes anuales distribuidos en las siete especialidades técnicas habilitadas (Construcción Civil, Contaduría Pública, Electrónica, Gastronomía, Mecánica Automotriz, Sistemas Informáticos y Textiles y Confecciones). De manera complementaria, la población integra al plantel docente, directivo y personal de secretaría académica encargado de la custodia documental y conformación de paralelos.

Dadas las características del contexto y la accesibilidad de los participantes en los periodos de prueba, se aplicó un muestreo no probabilístico por conveniencia intencional. La muestra humana final estuvo conformada por:

\begin{itemize}
    \item 60 estudiantes postulantes, seleccionados con criterio de representatividad entre las distintas unidades educativas de procedencia y opciones vocacionales, quienes participaron en la evaluación diagnóstica y en las pruebas de usabilidad de la nueva interfaz y del autoservicio de credenciales.
    \item 5 funcionarios institucionales (Director del módulo, responsable de secretaría académica y docentes de especialidad), quienes validaron el módulo de matriculación documental, la asignación automática de cursos y paralelos, y el panel de reportes oficiales.
\end{itemize}

\subsection{Unidades experimentales de prueba técnica}

Independientemente de los sujetos humanos, las unidades experimentales de prueba técnica corresponden a los ciclos operativos y réplicas de laboratorio del sistema V2.0 sometidos a ensayo controlado. Se ejecutaron réplicas de carga y auditoría transaccional sobre las siete especialidades: 5 corridas con Apache Benchmark 2.4.58 a $>$100 peticiones concurrentes por minuto, con latencia objetivo $<$1.2 s, y auditoría continua de logs del motor PostgreSQL 16.4 bajo bloqueo pesimista \texttt{FOR UPDATE}, verificando 0 sobrecupos y 100\% de consistencia de cupos en base de datos.
